\documentclass[11pt]{article}
\usepackage[margin=1in]{geometry}
\usepackage{amsmath,amssymb,amsthm,bm}
\usepackage{mathtools}
\usepackage{enumitem}
\usepackage[dvipsnames]{xcolor}
\usepackage{tcolorbox}
\tcbuselibrary{skins,breakable}
\usepackage{hyperref}
\hypersetup{colorlinks=true,linkcolor=NavyBlue,urlcolor=NavyBlue,
  pdftitle={The Source-Term Surrogate for Stiff Chemical Kinetics}}

\DeclareMathOperator{\diag}{diag}
\DeclareMathOperator{\spn}{span}
\DeclareMathOperator*{\argmin}{arg\,min}
\newcommand{\R}{\mathbb{R}}
\newcommand{\vv}[1]{\bm{#1}}
\newcommand{\Yv}{\vv{Y}}
\newcommand{\phiv}{\vv{\phi}}
\newcommand{\Sv}{\vv{S}}
\newcommand{\Rmap}{\mathcal{R}}
\newcommand{\Net}{\mathcal{N}}
\newcommand{\dt}{\Delta t}

\theoremstyle{plain}
\newtheorem{theorem}{Theorem}[section]
\newtheorem{proposition}[theorem]{Proposition}
\newtheorem{lemma}[theorem]{Lemma}
\theoremstyle{definition}
\newtheorem{definition}[theorem]{Definition}
\newtheorem{remark}[theorem]{Remark}

\newtcolorbox{engbox}[1]{breakable,enhanced,colback=NavyBlue!4,
  colframe=NavyBlue!55,fonttitle=\bfseries,title={#1}}

\title{The Source-Term Surrogate for Stiff Chemical Kinetics\\
{\large Why the flow map is learnable where the vector field is not}}
\author{}
\date{}

\begin{document}
\maketitle

\begin{abstract}
Physics-informed neural networks fail on stiff chemical kinetics for a reason
that is structural rather than incidental: the training residual inherits the
Lipschitz constant of the chemical source term, which for hydrocarbon and
hydrogen mechanisms spans eight to twelve decades. This note derives that
obstruction, then shows that the object a working surrogate actually learns ---
the finite-time \emph{flow map} of the chemistry substep --- is a contraction
precisely in the directions that make the vector field stiff. Stiffness, which
destroys the residual formulation, is what makes the flow-map formulation
well-posed. We then give the constraints that make such a surrogate usable
inside a CFD solver: exact element and mass conservation by linear projection,
temperature recovered from enthalpy rather than predicted, positivity, and
per-species scale separation. We close with an evaluation protocol designed
against the documented failure mode of this literature.
\end{abstract}

\section{Where the surrogate goes}

Consider a reacting flow with $n_s$ species. Writing $\rho$ for density,
$\vv{u}$ for velocity and $Y_k$ for the mass fraction of species $k$,
\begin{equation}
\frac{\partial (\rho Y_k)}{\partial t} + \nabla\cdot(\rho\vv{u}Y_k)
  \;=\; \nabla\cdot\!\left(\rho D_k \nabla Y_k\right) \;+\; \dot\omega_k,
\qquad k = 1,\dots,n_s,
\label{eq:species}
\end{equation}
where $\dot\omega_k$ is the chemical production rate. Essentially every
production solver advances \eqref{eq:species} by operator splitting. In Strang
form, over a step $\dt$,
\begin{equation}
\Phi_{\dt} \;=\; \mathcal{T}_{\dt/2}\circ\mathcal{C}_{\dt}\circ\mathcal{T}_{\dt/2},
\label{eq:strang}
\end{equation}
with $\mathcal{T}$ the transport (advection--diffusion) operator and
$\mathcal{C}$ the chemistry operator. The essential fact is that
$\mathcal{C}_{\dt}$ is \emph{spatially local}: in each cell it is an
initial-value problem in the thermochemical state alone,
\begin{equation}
\phiv \;=\; (T,\,Y_1,\dots,Y_{n_s})^{\!\top} \in \R^{n_s+1},
\qquad
\frac{d\phiv}{dt} \;=\; \Sv(\phiv),
\label{eq:ode}
\end{equation}
with no spatial coupling and no neighbours. The chemistry substep is therefore
the same $\R^{n_s+1}\to\R^{n_s+1}$ problem solved independently in every cell,
$10^6$ to $10^9$ times per simulation. In a detailed-chemistry LES it routinely
consumes $70$--$90\%$ of wall time.

\begin{definition}[Reaction map]
The \emph{reaction map} at step $\dt$ is the time-$\dt$ flow of \eqref{eq:ode},
\begin{equation}
\Rmap_{\dt}: \R^{n_s+1}\to\R^{n_s+1},
\qquad
\Rmap_{\dt}(\phiv_0) \;=\; \phiv_0 + \int_{0}^{\dt}\Sv\big(\phiv(\tau)\big)\,d\tau ,
\label{eq:reactionmap}
\end{equation}
where $\phiv(\cdot)$ solves \eqref{eq:ode} from $\phiv(0)=\phiv_0$.
\end{definition}

\begin{engbox}{What is being replaced}
The surrogate replaces \emph{one call} to the stiff ODE integrator --- CVODE,
Radau, or a specialised implicit method --- with one forward pass. It does not
replace the CFD solver, the transport operator, the pressure solve, or the
turbulence closure. Those remain exactly as they were, and they continue to
enforce the conservation laws they always enforced. This is the entire reason
the approach is tractable, and it is what distinguishes it from a
physics-informed network that attempts to represent the solution field itself.
\end{engbox}

\section{Why the residual formulation fails}

\subsection{The stiffness of the vector field}

Let $J(\phiv) = \partial\Sv/\partial\phiv \in \R^{(n_s+1)\times(n_s+1)}$ be the
Jacobian of the chemical source, with eigenvalues $\{\lambda_i\}$. Define the
stiffness ratio on a trajectory,
\begin{equation}
\varsigma \;=\;
\frac{\max_i \left|\operatorname{Re}\lambda_i\right|}
     {\min_{i:\,\operatorname{Re}\lambda_i\neq 0} \left|\operatorname{Re}\lambda_i\right|}.
\label{eq:stiffratio}
\end{equation}
For a hydrogen--air mechanism $\varsigma \sim 10^{8}$; for detailed
hydrocarbon mechanisms $\varsigma$ reaches $10^{12}$. The fast eigenvalues
belong to radical species --- $\mathrm{H}$, $\mathrm{O}$, $\mathrm{OH}$,
$\mathrm{HO_2}$ --- whose concentrations equilibrate on timescales of
$10^{-10}$ to $10^{-8}\,$s, while the slow modes governing heat release and
fuel consumption evolve on $10^{-4}$ to $10^{-2}\,$s.

\subsection{The obstruction}

A physics-informed network posits $\phiv_\theta(t)$ represented by a network
with parameters $\theta$, and minimises the residual
\begin{equation}
L(\theta) \;=\; \int_0^{\mathcal{T}}
  \left\| \frac{d\phiv_\theta}{dt}(t) - \Sv\big(\phiv_\theta(t)\big) \right\|^2 dt .
\label{eq:pinnloss}
\end{equation}

\begin{proposition}[Residual conditioning]
\label{prop:cond}
Let $g(\theta) = \tfrac{d\phiv_\theta}{dt} - \Sv(\phiv_\theta)$ denote the
residual and let $P_\theta = \partial\phiv_\theta/\partial\theta$ be the
parameter Jacobian. The Gauss--Newton approximation to the Hessian of
\eqref{eq:pinnloss} is
\begin{equation}
H \;\approx\; \int_0^{\mathcal{T}}
  P_\theta^{\!\top}\big(\partial_t - J\big)^{\!\top}\big(\partial_t - J\big) P_\theta \, dt .
\end{equation}
The spectrum of $(\partial_t - J)^{\!\top}(\partial_t - J)$ contains the squared
eigenvalue magnitudes $|\lambda_i|^2$. Hence, absent preconditioning,
\begin{equation}
\kappa(H) \;\gtrsim\; \varsigma^2 ,
\end{equation}
and first-order optimisation converges at a rate governed by
$1 - \mathcal{O}(\kappa(H)^{-1})$.
\end{proposition}

For $\varsigma = 10^8$ this gives $\kappa \gtrsim 10^{16}$, which is the
reciprocal of double-precision epsilon. No learning-rate schedule, adaptive
optimiser, or loss reweighting repairs a condition number at that magnitude;
the fast modes saturate the residual and the slow modes --- the ones carrying
the heat release one actually wants --- contribute gradient components below
numerical noise.

\begin{engbox}{The same obstruction, in classical language}
This is not a machine-learning pathology. It is the identical obstruction that
forbids explicit time integration of stiff systems. An explicit method requires
$\dt \lesssim 2/\max_i|\lambda_i|$, which for $\varsigma=10^{8}$ means resolving
$10^{-10}\,$s steps to follow $10^{-2}\,$s dynamics. The classical response was
not a better explicit method: it was to change formulation, to implicit schemes
that solve a linear system with $J$ rather than stepping along it. Gradient
descent on \eqref{eq:pinnloss} is an explicit method in disguise, and it fails
for the same reason.
\end{engbox}

\section{Why the flow map is learnable}

The surrogate does not learn $\Sv$. It learns $\Rmap_{\dt}$. This is a
different object with different conditioning, and the difference is the whole
argument.

\begin{proposition}[The flow map contracts the stiff directions]
\label{prop:contract}
Let $\Sv$ be linearised about an equilibrium $\phiv^\ast$ on the slow manifold,
$\Sv(\phiv)\approx J(\phiv-\phiv^\ast)$ with $J$ diagonalisable,
$J = V\Lambda V^{-1}$, $\Lambda=\diag(\lambda_1,\dots,\lambda_{n_s+1})$ and
$\operatorname{Re}\lambda_i \le 0$. Then
\begin{equation}
\Rmap_{\dt}(\phiv) - \phiv^\ast
  \;=\; V \exp(\Lambda \dt) V^{-1} (\phiv - \phiv^\ast),
\end{equation}
so the singular values of $D\Rmap_{\dt}$ in the eigenbasis are
$e^{\operatorname{Re}\lambda_i \dt}$. Consequently
\begin{equation}
\operatorname{Lip}(\Rmap_{\dt}) \;\le\; \kappa(V)\,
  e^{\dt\,\max_i \operatorname{Re}\lambda_i} \;\le\; \kappa(V),
\end{equation}
whereas $\operatorname{Lip}(\Sv) = \|J\| \sim \max_i|\lambda_i|$.
\end{proposition}

\begin{proof}
Immediate from $\Rmap_{\dt} = \exp(J\dt)$ in the linearised setting and the
spectral mapping theorem. The bound on the Lipschitz constant follows since
$\operatorname{Re}\lambda_i\le 0$ gives $e^{\operatorname{Re}\lambda_i\dt}\le 1$.
\end{proof}

The content of Proposition~\ref{prop:contract} is worth stating plainly. In the
fast directions, $|\lambda_i|\dt \gg 1$, so $e^{\lambda_i \dt}\approx 0$: the
flow map \emph{annihilates} exactly those directions in which the vector field
is violent. What survives after $\dt$ is motion on the slow manifold, whose
dimension is typically $5$ to $15$ regardless of $n_s$. The map is smooth,
bounded, non-expansive, and effectively low-dimensional.

\begin{engbox}{The inversion}
Stiffness is the obstruction for the residual formulation and the \emph{gift}
for the flow-map formulation. A large $\varsigma$ means the fast transients have
already decayed within one $\dt$, so the function being approximated is a
projection onto a low-dimensional attracting manifold. The stiffer the
chemistry, the more the flow map contracts, and the easier it is to learn. This
is why the same problem that defeats a PINN yields to a supervised surrogate,
and it is not a matter of engineering effort or network capacity.
\end{engbox}

\begin{remark}
This is also why the idea long predates machine learning. In-situ adaptive
tabulation and flamelet-manifold methods are, in this language, nearest-neighbour
and low-rank approximations of the same $\Rmap_{\dt}$. A neural surrogate is a
better interpolant for the same object, not a new physical idea.
\end{remark}

\section{Constructing the surrogate}

\subsection{Target and scale separation}

Predict the \emph{increment}, not the state:
\begin{equation}
\Delta\phiv \;=\; \Rmap_{\dt}(\phiv_0) - \phiv_0 .
\end{equation}
Mass fractions span more than twenty decades, from $Y \sim 1$ for the bath gas
to $Y\sim10^{-20}$ for trace radicals, and a uniformly-weighted $\ell^2$ loss on
raw $Y$ is dominated entirely by $\mathrm{N_2}$. Two transformations are used,
and they matter more than architecture.

\paragraph{Logarithmic (or Box--Cox) transform.} With a floor $\epsilon>0$,
\begin{equation}
\psi_k \;=\; \ln\!\left(Y_k + \epsilon\right),
\qquad\text{or}\qquad
\psi_k \;=\; \frac{(Y_k+\epsilon)^{\nu}-1}{\nu},
\label{eq:boxcox}
\end{equation}
the latter reducing to the former as $\nu\to0$. This makes relative error the
quantity being minimised, which is the quantity that matters for a radical
species four decades below the fuel.

\paragraph{Per-species standardisation.} After transformation,
\begin{equation}
\tilde\psi_k \;=\; \frac{\psi_k - \mu_k}{\sigma_k},
\qquad
\mu_k=\mathbb{E}[\psi_k],\quad \sigma_k^2=\operatorname{Var}[\psi_k],
\end{equation}
with the statistics computed over the training set. Each species then
contributes comparably to the loss.

\subsection{Exact conservation by projection}

Let $E\in\R^{n_e\times n_s}$ have entries
\begin{equation}
E_{jk} \;=\; a_{jk}\,\frac{W_j}{W_k},
\end{equation}
where $a_{jk}$ is the number of atoms of element $j$ in species $k$, and $W$
denotes molar mass. Element conservation during the chemistry substep is the
\emph{linear} statement
\begin{equation}
E\,\Delta\Yv \;=\; \vv{0},
\label{eq:elemcons}
\end{equation}
and total mass conservation $\sum_k \Delta Y_k = 0$ is the additional row
$\vv{1}^{\!\top}\Delta\Yv=0$. Append it to $E$. Rather than penalising
\eqref{eq:elemcons} in the loss, impose it exactly: let $\Net_\theta$ be the
raw network output and project,
\begin{equation}
\boxed{\;
\Delta\Yv \;=\; \left(I - E^{+}E\right)\Net_\theta(\phiv_0),
\qquad E^{+}=E^{\!\top}(EE^{\!\top})^{-1}.
\;}
\label{eq:projection}
\end{equation}
Since $\left(I-E^{+}E\right)$ is the orthogonal projector onto $\ker E$, the
constraint holds to machine precision for every input, trained or not. The cost
is one $n_e\times n_s$ matrix product with $n_e\le5$, negligible against the
forward pass, and $E^{+}$ is precomputed once.

\begin{engbox}{Hard constraints beat penalties here}
A conservation penalty $\lambda\|E\Delta\Yv\|^2$ trades conservation against
accuracy through a weight nobody can choose principledly, and it is violated on
exactly the out-of-distribution states where violation is most damaging.
Projection \eqref{eq:projection} costs less, cannot be tuned wrong, and holds
under distribution shift. Where a constraint is linear, project; reserve
penalties for constraints that are not.
\end{engbox}

\subsection{Temperature from enthalpy, not from the network}

At constant pressure the chemistry substep conserves specific enthalpy:
\begin{equation}
h \;=\; \sum_{k=1}^{n_s} Y_k h_k(T) \;=\; \text{const},
\qquad
h_k(T) = h_k^{\circ} + \int_{T^{\circ}}^{T} c_{p,k}(T')\,dT'.
\label{eq:enthalpy}
\end{equation}
Do not give the network a temperature output. Predict $\Delta\Yv$ only, then
recover $T$ by solving the scalar equation \eqref{eq:enthalpy} for $T$ with
Newton's method,
\begin{equation}
T^{(m+1)} \;=\; T^{(m)} -
  \frac{\sum_k Y_k h_k(T^{(m)}) - h}{\sum_k Y_k c_{p,k}(T^{(m)})},
\end{equation}
which converges in two or three iterations from the previous temperature. This
reduces the output dimension by one, removes a large-magnitude output that would
otherwise dominate the loss, and makes energy conservation exact rather than
approximate.

\subsection{Positivity and the simplex}

Mass fractions must satisfy $Y_k\ge0$ and $\sum_k Y_k=1$. Predicting in the
transformed variable \eqref{eq:boxcox} and inverting gives $Y_k>0$
automatically. Renormalise afterwards,
\begin{equation}
Y_k \leftarrow \frac{\max(Y_k,0)}{\sum_j \max(Y_j,0)} ,
\end{equation}
applied \emph{after} the projection \eqref{eq:projection} and accepted as a
final $\mathcal{O}(\text{machine }\epsilon)$ correction.

\section{The failure mode that matters: a posteriori drift}

A surrogate with excellent single-step error can still destroy a simulation.
Writing $\hat\Rmap$ for the surrogate and $\varepsilon_n$ for the error after
$n$ applications, one step of error propagation gives
\begin{equation}
\varepsilon_{n+1} \;\le\; \operatorname{Lip}(\Rmap_{\dt})\,\varepsilon_n
   + \left\|\hat\Rmap - \Rmap_{\dt}\right\|_{\infty}.
\label{eq:accum}
\end{equation}
By Proposition~\ref{prop:contract}, $\operatorname{Lip}(\Rmap_{\dt})\le1$ on the
stiff directions, so error does \emph{not} amplify along the attracting
manifold --- which is the structural reason the approach is stable at all. The
danger is elsewhere: the coupled system revisits states the training set never
covered, because transport carries cells into composition space that
one-dimensional flame data does not populate. Two mitigations are standard.

\begin{enumerate}[leftmargin=*]
\item \textbf{Manifold-aware sampling.} Draw training states from canonical
      flames --- freely propagating premixed, counterflow diffusion,
      homogeneous ignition --- which populate the attracting manifold densely.
\item \textbf{Physics-aware augmentation.} Interpolate between neighbouring
      sampled states and add constrained random perturbations that respect
      \eqref{eq:elemcons} and \eqref{eq:enthalpy}, thickening the training
      distribution transverse to the manifold without leaving the physically
      admissible set.
\end{enumerate}

\section{Evaluation protocol}

The published record in this area carries a documented reporting problem: of
articles claiming that a machine-learned method outperforms a standard
numerical method on a fluid-related PDE, $79\%$ ($60$ of $76$) compared against
a weak baseline, with outcome-reporting and publication bias both present.\!%
\footnote{N.~McGreivy and A.~Hakim, \emph{Weak baselines and reporting biases
lead to overoptimism in machine learning for fluid-related partial differential
equations}, Nature Machine Intelligence \textbf{6}, 1256--1269 (2024).}
The protocol below is written against that finding.

\begin{enumerate}[leftmargin=*]
\item \textbf{The baseline is a tuned CVODE at the tolerance actually required},
      not at a tolerance chosen to make it slow. Report the tolerance.
\item \textbf{Include the real incumbent.} The production alternative is not
      naive direct integration; it is tabulation --- ISAT or a flamelet
      manifold. A surrogate that beats CVODE but loses to ISAT has not earned
      its complexity.
\item \textbf{Report a posteriori, not a priori.} Single-step $L^2$ error on a
      held-out set is not the claim. The claim is a coupled CFD run: ignition
      delay, laminar flame speed, extinction strain rate, and the full
      temperature and major-species profiles against the direct-integration
      reference.
\item \textbf{Report end-to-end wall time}, including the projection, the
      Newton solve for $T$, and the host--device transfers --- not kernel time.
\item \textbf{State the conservation residuals} $\|E\Delta\Yv\|_\infty$ and the
      enthalpy drift over the whole run.
\item \textbf{Report where it breaks.} Name the equivalence ratio, pressure and
      temperature range outside which the surrogate departs from the reference,
      and publish that boundary with the model.
\end{enumerate}

\begin{engbox}{The honest target}
Reported gains for this formulation are $\mathcal{O}(10^2)$ on isolated
chemistry evaluation against CVODE and roughly $10$--$20\times$ end-to-end in
coupled CFD. The gap between those two numbers is the point: once chemistry
stops dominating, Amdahl's law caps the benefit, and further speedup must come
from transport. A surrogate is worth building when a design sweep or an
uncertainty quantification study needs $10^5$--$10^6$ runs and tolerates
percent-level error. It is not worth building to make a single high-fidelity
run faster.
\end{engbox}

\end{document}
